\chapter{Auxiliary theory}
\label{chap:auxiliary}

The Galois-theoretic part rests on general facts about multiquadratic extensions, invariants
of~$2$-groups, wreath products and even polynomials. They are formalized in Mathlib generality under
\texttt{QuadraticIterates/Mathlib/}, except for the iteration sequences of even polynomials
(\texttt{QuadraticIterates/EvenPoly.lean}).

\section{The iteration sequence of an even polynomial}

\begin{definition}[Even polynomials and their iteration sequences]
  \label{def:gammaseq}
  \lean{QuadraticIterates.EvenPoly, QuadraticIterates.evenPoly_iff_comp_neg_X,
    QuadraticIterates.EvenPoly.eval_neg, QuadraticIterates.EvenPoly.eval_congr,
    QuadraticIterates.EvenPoly.map, QuadraticIterates.gammaSeq, QuadraticIterates.gammaSeq_zero,
    QuadraticIterates.gammaSeq_one, QuadraticIterates.gammaSeq_succ,
    QuadraticIterates.map_gammaSeq, QuadraticIterates.gammaSeq_eq_ite_even_of_add_two_eq,
    QuadraticIterates.gammaSeq_add, QuadraticIterates.intCast_gammaSeq,
    QuadraticIterates.EvenPoly.dvd_eval_sub, QuadraticIterates.EvenPoly.iterate_eval_neg}
  \leanok
  A polynomial~$g$ is \emph{even} if $g \in R[X^2]$. Its iteration sequence with
  sign~$\varepsilon$ is $\gamma_1 = \varepsilon g(0)$, $\gamma_{n+1} = g(\gamma_n)$ for $n \ge 1$;
  it commutes with ring homomorphisms, and if $\gamma_{n_0 + 2} = \gamma_{n_0}$ for some $n_0 \ge
  1$, it is~$2$-periodic from~$n_0$ on.
  \formalnote{The sequence~$c$ is the case $g = X^2 + a$, $\varepsilon = -1$.}
\end{definition}

\section{Multiquadratic extensions}

\begin{lemma}[Degree of a multiquadratic extension]
  \label{lem:multiquadratic}
  \lean{multiquadraticRelations, mem_multiquadraticRelations, finrank_adjoin_image_eq_two_pow_sub,
    finrank_adjoin_range_eq_two_pow_sub, IntermediateField.finrank_adjoin_insert,
    IntermediateField.finrank_adjoin_insert_of_not_isSquare,
    IntermediateField.finrank_adjoin_sqrt_le, IntermediateField.finrank_adjoin_finset_sqrt_le,
    IntermediateField.finrank_adjoin_sqrt_le_two_pow_ncard,
    IntermediateField.finrank_adjoin_sqrt_eq, IntermediateField.finrank_adjoin_sqrt_eq_one_iff,
    finrank_multiquadraticRelations_le,
    finrank_multiquadraticRelations_insert}
  \leanok
  Let~$F$ be a field of characteristic $\ne 2$ and $(x_i)_{i \in s}$ a finite family of square
  roots of elements $r_i \in F^*$ inside an extension. Then
  $[F(x_i : i \in s) : F] = 2^{\#s - \dim_{\FF_2} R(s)}$, where $R(s) \subseteq \FF_2^{s}$ is the
  space of \emph{relations}: the~$\varepsilon$ with $\prod_{i : \varepsilon_i = 1} r_i$ a square
  in~$F$. (The family need not be injective: a repeated root counts once in~$\#s$ and once
  in~$R(s)$.) In particular $[F(x_i) : F] \le 2^{\#s}$, with equality iff there are no relations;
  and adjoining one further square root either doubles the degree or leaves it unchanged.
\end{lemma}
\begin{proof}
  \leanok
  \uses{lem:descent}
  Induction on~$\#s$, the step being the descent Lemma~\ref{lem:descent}: the degree doubles
  exactly when the new radicand is not a square in the field generated so far.
\end{proof}

\begin{lemma}[Square descent]
  \label{lem:descent}
  \uses{def:twoindep}
  \lean{IntermediateField.adjoin_insert, IntermediateField.exists_mem_adjoin_sq_eq_algebraMap_iff,
    IntermediateField.exists_mem_adjoin_sq_eq_algebraMap_of_isSquare_mul_prod,
    IntermediateField.exists_mem_adjoin_simple_sq_eq_algebraMap_iff,
    IntermediateField.isSquare_algebraMap_iff,
    IntermediateField.not_isSquare_algebraMap_of_sqrt_notMem,
    IntermediateField.isSquare_algebraMap_bot_iff, IntermediateField.isSquare_algebraMap_iff_mem,
    IntermediateField.mem_adjoin_sqrt_iff,
    sqClass_mul, sqClass_prod, sqClass_prod_zpow, sqClass_eq_zero_iff, sqClass_pow, sqClass_mul_sq,
    sqClass_div_sq, ZModModule.zsmul_eq_natAbs_nsmul, isSquare_prod_iff_sum_sqClass_eq_zero}
  \leanok
  An element of~$F^*$ is a square in a multiquadratic extension $F(\sqrt{c_y} : y \in s)$ if and
  only if it becomes a square in~$F$ after multiplication by a subproduct of the~$c_y$.
\end{lemma}
\begin{proof}
  \leanok
  Adjoining one square root at a time; in each step, an element of the smaller field whose square
  root lies in the quadratic extension differs from a square by the radicand.
\end{proof}

\begin{lemma}[Relations among the shifted roots]
  \label{lem:relations}
  \uses{lem:multiquadratic, def:twoindep}
  \lean{rootRelations, mem_rootRelations, rootRelations_eq_bot_iff, finrank_rootRelations_le,
    multiquadraticRelations_univ, finrank_multiquadraticRelations,
    one_mem_rootRelations_iff, Function.ExtendByZero.linearMap_injective,
    exists_mem_multiquadraticRelations_insert_iff, multiquadraticRelations_insert_inf_ker_proj,
    indicator_mem_multiquadraticRelations_iff}
  \leanok
  For a family $(r_i)_{i \in \iota}$ of nonzero elements of~$F$, the relation space
  $R \subseteq \FF_2^{\iota}$ of the full family is the space of relations of
  Lemma~\ref{lem:multiquadratic} for $s = \iota$; it vanishes iff the classes of the~$r_i$ in
  $F^*/F^{*2}$ are linearly independent, and $\dim R \le \#\iota$. The all-ones vector lies
  in~$R$ if and only if~$\prod_i r_i$ is a square in~$F$.
\end{lemma}
\begin{proof}
  \uses{def:twoindep, lem:descent}
  \leanok
  Unfolding the definition of the relation space.
\end{proof}

\begin{lemma}[Invariants of a~$p$-group]
  \label{lem:invariant}
  \uses{lem:relations}
  \lean{IsPGroup.exists_ne_zero_mem_fixedPoints_of_smul_mem,
    IsPGroup.one_mem_of_comp_smul_mem_of_ne_bot, comp_smul_mem_rootRelations, one_mem_rootRelations_of_ne_bot}
  \leanok
  Let a~$p$-group~$G$ act on an~$\FF_p$-vector space~$M$, and let $V \subseteq M$ be a finite
  nonzero~$G$-stable subspace. Then $V^G \ne 0$. If moreover $M = \FF_p^{\iota}$ and~$G$ acts by
  permuting the coordinates \emph{transitively}, then the all-ones vector lies in~$V$.
  (We use this with $p = 2$.)
\end{lemma}
\begin{proof}
  \leanok
  For the first statement, $\#V$ is a power of~$p$, and $\#V \equiv \#V^G \bmod p$ by Mathlib's
  \lname{IsPGroup.card\_modEq\_card\_fixedPoints}; since $0 \in V^G$, this gives $\#V^G > 1$
  (Mathlib's \lname{IsPGroup.exists\_fixed\_point\_of\_prime\_dvd\_card\_of\_fixed\_point}).
  For the second, a nonzero invariant vector has constant coordinates along the orbits of~$G$
  on~$\iota$, and there is only one orbit; scaling makes the constant equal to~$1$.
\end{proof}

\section{Wreath products}

\begin{lemma}[Sylow~$2$-subgroups of symmetric groups]
  \label{lem:sylowwreath}
  \lean{Sylow.mulEquivIteratedWreathProduct, IteratedWreathProduct.card,
    IsPGroup.exists_injective_monoidHom_iteratedWreathProduct,
    MonoidHom.nonempty_mulEquiv_iff_card_eq}
  \leanok
  A Sylow~$2$-subgroup of the symmetric group on~$2^n$ letters is isomorphic to~$\Ct{n}$, which
  has order $2^{2^n - 1}$; hence a~$2$-group of permutations of at most~$2^n$ points embeds
  into~$\Ct{n}$. Moreover, an injective homomorphism between finite groups of equal order is an
  isomorphism, which is how ``$\Om_n \cong \Ct{n}$'' is obtained from an embedding plus a degree
  count.
\end{lemma}
\begin{proof}
  \leanok
  The first statement is in Mathlib (\texttt{Sylow.mulEquivIteratedWreathProduct}); the second
  follows by extending the action to~$2^n$ points and passing to a Sylow subgroup, and the last
  one is elementary.
\end{proof}

\begin{lemma}[The abelian quotient of~$\Ct{n}$]
  \label{lem:wreathab}
  \lean{IteratedWreathProduct.card_monoidHom_multiplicative_zmod_two,
    IteratedWreathProduct.card_monoidHom, RegularWreathProduct.card_monoidHom,
    RegularWreathProduct.monoidHomEquiv, Nat.card_monoidHom_multiplicative_zmod_two,
    Nat.card_monoidHom_multiplicative_zmod, Nat.card_addMonoidHom_zmod,
    IntermediateField.isGalois_adjoin_of_sq_eq_algebraMap,
    IntermediateField.algEquiv_adjoin_sq_eq_one, AlgEquiv.apply_eq_or_eq_neg_of_sq_eq_algebraMap,
    IntermediateField.finrank_adjoin_le_card_monoidHom,
    IntermediateField.mem_adjoin_of_sq_eq_algebraMap_of_card_le}
  \leanok
  $\#\Hom(\Ct{n}, C_2) = 2^n$; equivalently, the largest elementary abelian~$2$-quotient
  of~$\Ct{n}$ is~$C_2^n$. A field extension generated by square roots of base-field elements is
  Galois with elementary abelian group of exponent~$2$, so its degree equals the number of
  homomorphisms of its Galois group to~$C_2$; inside a finite Galois extension~$E/L$ such a
  subfield therefore has degree at most $\#\Hom(\Gal(E/L), C_2)$, and one of that maximal degree
  contains every square root of an element of~$L$ lying in~$E$.
\end{lemma}
\begin{proof}
  \leanok
  \uses{lem:descent, lem:multiquadratic}
  For a regular wreath product~$D \wr Q$ and an abelian group~$A$, a homomorphism $D \wr Q \to A$
  is the same as a pair of homomorphisms $D \to A$ and $Q \to A$ (on the base group~$D^Q$ it
  takes the same values on all coordinates, since these are conjugate under~$Q$), giving
  $\#\Hom(D \wr Q, A) = \#\Hom(D, A) \cdot \#\Hom(Q, A)$; induction on~$n$ then gives
  $\#\Hom(\Ct{n}, C_2) = \#\Hom(C_2, C_2)^n = 2^n$. The statement about elementary abelian
  groups is the identity $\#G = \#\Hom(G, C_p)$ for a finite~$\FF_p$-vector space~$G$, i.e.\
  $\dim G^{\vee} = \dim G$ (for $p = 2$, commutativity follows from the exponent). Finally, if a
  subfield~$M$ of~$E$ generated by square roots has the maximal degree and $x \in E \setminus M$
  with $x^2 \in L$, then~$x^2$ is not a square in~$M$ (Lemma~\ref{lem:descent}), so
  $[M(x) : L] = 2[M : L]$ (Lemma~\ref{lem:multiquadratic}), which is impossible.
\end{proof}

\section{Even polynomials}

\begin{lemma}[Even polynomials are polynomials in~$X^2$]
  \label{lem:evenpoly}
  \lean{Polynomial.exists_eq_comp_X_sq_add_C_of_comp_neg_X_eq,
    Polynomial.eq_expand_two_contract_of_comp_neg_X_eq, Polynomial.expand_two_comp_neg_X,
    Polynomial.X_sq_add_C_comp_neg_X, Polynomial.comp_neg_X_eq_or_eq_neg_of_associated,
    Polynomial.natDegree_comp_neg_X, Polynomial.expand_contract_of_forall_coeff_eq_zero,
    Polynomial.coeff_eq_zero_of_comp_neg_X_eq_of_not_two_dvd, Polynomial.comp_dvd_comp_iff}
  \leanok
  Over a domain of characteristic $\ne 2$, a polynomial fixed by $X \mapsto -X$ is a polynomial
  in~$X^2$, hence in~$X^2 + c$ for any~$c$. A polynomial merely \emph{associated} to its
  reflection satisfies $p(-X) = p$ or $p(-X) = -p$. Over a field, $p \circ r \dvd q \circ r$
  implies~$p \dvd q$ for every nonconstant~$r$.
\end{lemma}
\begin{proof}
  \leanok
  The odd coefficients vanish because $2 \ne 0$; for the second statement, the unit relating~$p$
  to $p(-X)$ is a constant squaring to~$1$. For the last statement, the remainder of~$q$
  modulo~$p$, composed with~$r$, is divisible by~$p \circ r$ but has smaller degree, so it vanishes.
\end{proof}

\begin{lemma}[Pairing the irreducible factors]
  \label{lem:involution}
  \lean{Multiset.exists_add_map_of_involutive, normalizedFactors_map_mulEquiv_eq,
    normalize_map_normalize, exists_normalizedFactors_eq_add_map}
  \leanok
  A fixed-point-free involution of a multiset splits it as $N + \tau(N)$. Consequently,
  if~$\sigma$ is an involutive automorphism of a UFD and~$p$ is a reducible element with
  $\sigma(p) \sim p$ whose~$\sigma$-invariant divisors (up to associates) are all trivial, then
  the normalized factors of~$p$ pair up under $\mathrm{normalize} \circ \sigma$. For
  $\sigma = (X \mapsto -X)$ on the polynomial ring over a field, this yields the factorization
  $\pm p = g \cdot g(-X)$ of an even polynomial without nontrivial even divisors.
\end{lemma}
\begin{proof}
  \leanok
  Induction on the multiset, removing an orbit $\{q, \tau q\}$ at a time; that
  $\mathrm{normalize} \circ \sigma$ preserves the multiset of normalized factors is
  functoriality of factorization under~$\sigma$, and it has no fixed point there because a fixed
  normalized factor would be a nontrivial~$\sigma$-invariant divisor.
\end{proof}
